\scanpage{127}
\deleted{On a envie de définir des torseurs}
\[
\deleted{D_{d,d'}^*=D_{d,d'}\mid D_d^*\quad\text{ouvert de }D_{d,d'}}
\]
\deleted{(quand on fait l'\unsure{hyp}}

On suppose maintenant que la situation envisagée provient d'un diviseur à croisements normaux $D$ sur $X=D_0$, en déployant celui-ci. Alors \inserted{$\forall\,d\in\mathbb N$} la structure \inserted{de} déplo\unsure{y|i}ement induite sur les $D_d$ (en posant $D_{\deleted{\unsure{d}}i}^{\langle d\rangle}=D_{d+i,d}$ \mbox{-\kern.3em-\kern.3em-\kern.3em-} \unsure{p|je} me \unsure{dispenser|dispense} de développer au long et au large) avec les $=$ propriétés --- elle provient d'un diviseur à croisements normaux \inserted{$\Theta_d$ \deleted{\unsure{$(D(d))$|(cf)}}} sur $D_d$, dont le complémentaire du support est noté $D_d^*$.

\deleted{On note}
\[
\deleted{(203)\qquad D_{d,d'}^*=D_{d,d'}\mid D_d^*}
\]
\margin{\deleted{ouvert de $D_{d,d'}$ (complémentaire d'un diviseur à croisements normaux $\Delta(d,d')$ sur $D_{d,d'}$\unsure{)}}\par \deleted{(image inverse de $\unsure{\Delta(d)|\Theta(d)|D(d)}$}}

\deleted{On \deleted{\unsure{posera|pose} alors} \inserted{veut définir} \deleted{des} morphismes canoniques de groupoïdes}
\[
\deleted{(203)\qquad\boxed{\Pi_1P_d^*\xrightarrow{\unsure{\gamma|r}_{d,d'}}\Pi_{\unsure{1}}P_{d'}^*}\qquad\text{pour }d'<d}
\]
\deleted{``h\uuline{om de générisation}''}
\margin{126}
